
\section{Problem Definition} \label{sec:pro_def}

We study the optimization of MAS within the context of multi-step collaboration process. The definition of agents and the fundamental collaboration workflow are as follows:


\noindent \textbf{Definition 1: Agents.} An LLM-powered agent is governed by natural language prompts to generate solutions for completing a given task. Formally, we define an agent $a$ as  a function: $\mathcal{A}: \mathcal{P} \times \mathcal{I} \rightarrow \mathcal{O}$. Here, $p \in \mathcal{P}$ represents the instruction prompt defining the role and functionality of the agent. In an in-context learning setting, the prompt may also include few-shot examples serving as demonstrations to guide the agent’s behavior. The input $i \in \mathcal{I}$ typically comprises the query information, task description, and potentially instructions or solutions from other agents in MAS. The output $o \in \mathcal{O}$ generally encompasses the agent's generated analyses, reasoning, and solutions to the given task.


\noindent \textbf{Definition 2: Multi-Step Collaboration.} In this work, we consider multiple agents collaborating within a multi-step procedure. Typically, each step $s_t \in \mathcal{S}$ involves a subset of agents from the overall team, denoted as $\{a_{t,1}, a_{t,2}, ..., a_{t,n}\} \subseteq s_t$. This multi-step approach aims to enhance problem-solving by enabling agents at each step to leverage solutions and outcomes produced by agents in preceding steps. For instance, when solving a complex mathematical problem, agents in the initial step may first decompose the problem into a series of simpler subproblems. Agents in subsequent steps can then address each subproblem by integrating analyses and solutions generated earlier with relevant contextual information. 
Especially, we propose determining and refining this collaboration structure (e.g., selecting agents to participate at each step) using \textbf{LLM-based controllers}. The functionalities of these controllers are detailed in Section \ref{sec:five_dim}.


\endinput


We formulate our problem as a Multi-Objective Markov Decision Process (MOMDP) \cite{roijers2013survey}, which can be described by the tuple $<\mathcal{S}, \mathcal{A}, T, R,\mathcal{S}, \rho>$ given a specific environment that the agent can interact with. The difference between MOMDP and classic Markov Decision Process (MDP) it that the reward function $R(\bs_t, a_t)$ in MOMDP is represented by a multi-dimensional vector in which each dimension describes the reward for a specific objective. In contrast, reward $R$ in a standard MDP is always a scalar. Especially, in our problem setting for sequential recommendation, the environment is taken as all the users on the platform and the agent is the recommender system. A state $\bs_t \in \mathcal{S}$ is taken as the representation of users' preferences at a given timestep $t$. 
%Following previous research \cite{xin2022rethinking, yang2024generate, kang2018self}, we obtain $\bs_t$ by using a sequential model to encode user's historical interaction behaviours. 
The action space $\mathcal{A}$ includes all candidate items to be recommended, such that $|\cA|=|\cV|=N$ where $\cV$ denotes the set of all items. $T$ is the dynamic function to define the state transition probability. $R$ is the reward function, where $\br_t=R(s_t, a_t)$ means the agent receives a reward 
%\footnote{We assume a deterministic reward function to simplify our exposition following previous research.} 
$\br_t$ after taking an action $a_t$ under state $\bs_t$. $\rho$ is the distribution of the initial states.


% \renewcommand{\algorithmicrequire}{\textbf{Input:}}  \raisebox{0.5ex}{$\boldsymbol{\chi}$} 
% \renewcommand{\algorithmicensure}{\textbf{Output:}}
% \algnewcommand{\algorithmicand}{\textbf{ and }}
% \algnewcommand{\algorithmicor}{\textbf{ or }}
% \algnewcommand{\OR}{\algorithmicor}
% \algnewcommand{\AND}{\algorithmicand}
% \setlength{\textfloatsep}{1.5pt}
% \begin{algorithm}[t]  % algorithm
% 	\caption{\textbf{The VPMCR Scenario}}			% title
% 	\begin{algorithmic}[1]	% there is raw labels for each raw
% 	\Require		
% 		user $u$,
%         Target item set $\mathcal{V}_{target}$,
%         all attribute types $\mathcal{C}$,
%         clear preference space $\mathcal{C}_{CI}$,
%         clear preference space $\mathcal{C}_{VI}$,
%         all fine-grained attributes $\mathcal{P}$,
%         all items $\mathcal{V}$,
%         %The set of attributes preferred by the user $\mathcal{P}_{u}$,
%         the number of items to recommend $k$,
%         the maximum number of turns $T$;
% 	\Ensure recommendation result: success or fail;
%     \State User $u$ specifies an common attribute $p_{0}$ from ; 
%     \State Update:  $\mathcal{P}_{u}=\{p_{0}\}$;
%     $\mathcal{P}_{cand}=\mathcal{P} \setminus p_{0}$;
%     $\mathcal{V}_{cand}=\mathcal{V}_{p_{0}}$
%     \For{turn $t=1,2,3...T$}
%         \State Select an action $a$
%         \If{$a == a_{ask}$}
%             \State Select the top attribute $p$ from $\mathcal{P}_{cand}$
%             \If{$u$ accepts $p_{t}$}
%                 \State Update: $\mathcal{P}_{u}=\mathcal{P}_{u} \cup p$;  $\mathcal{V}_{cand} =  \mathcal{V}_{cand} \cap \mathcal{V}_{p}$
%             \EndIf
%                 \State Update: $\mathcal{P}_{cand} = \mathcal{P}_{cand} \setminus p$
            
%         \Else[{$a == a_{rec}$}]
%             \State Select the top-$k$ items $\mathcal{V}_{k}$ from $\mathcal{V}_{cand}$
%             \If{User accepts $\mathcal{V}_{k}$}
%                 \State \emph{ Recommendation succeeds}; Exit.
%             \Else[User rejects $\mathcal{V}_{k}$]
%                 \State Update: $\mathcal{V}_{cand} = \mathcal{V}_{cand} \setminus \mathcal{V}_{k}$
%             \EndIf
%         \EndIf
%     \EndFor
%     \State \emph{Recommendation fails}; Exit. 
% 	\end{algorithmic}
% 	\label{algo:MCR}
% \end{algorithm}

% The conversation state representation at each turn is represented by $s_t$, the action taken by the agent is $a_t$, and the reward received for taking action $a_t$ in state $s_t$ is $r(s_t, a_t)$. The target policy $\pi^{*}$: $\arg \max _{\pi \in \Pi} \mathbb{E}\left[\sum_{t=0}^{T} r\left(s_{t}, a_{t}\right)\right]$. 




